% ============================================================================
%  Chapter 4 --- System Architecture
% ============================================================================
\chapter{System Architecture}
\label{ch:architecture}

\section{Introduction}

This chapter presents the architecture of the AI Operations Agent platform. We describe the system at four levels of abstraction: the high-level CEM--CVM positioning (Section~\ref{sec:arch-highlevel}), the container-level service decomposition (Section~\ref{sec:arch-containers}), the data flow through the pipeline (Section~\ref{sec:arch-pipeline}), and the data lake architecture (Section~\ref{sec:arch-datalake}).

\section{High-Level Architecture: CEM $\to$ AI Agent $\to$ CVM}
\label{sec:arch-highlevel}

At the highest level, the system is the intelligence bridge between two Huawei ecosystem components:

\begin{itemize}
  \item \textbf{Upstream (CEM/SmartCare):} Provides KPI/KQI/CEI scores, demarcation results, and experience alerts. In our implementation, this data arrives as real OSS exports from Tunisie Telecom (CSV/Excel/JSON files).
  \item \textbf{Downstream (CVM):} Consumes churn predictions, SLA risk scores, anomaly alerts, upsell triggers, and revenue impact correlations. In our implementation, CVM reads these from our REST API.
\end{itemize}

The AI Agent sits in the middle and performs four core functions:
\begin{enumerate}
  \item \textbf{Ingestion:} Reads OSS and BSS data, anonymises sensitive fields, maps to the internal schema.
  \item \textbf{Analysis:} Runs ML inference --- SLA risk scoring, network anomaly detection, revenue anomaly detection.
  \item \textbf{Correlation:} Computes statistical correlations between OSS degradation and BSS impact.
  \item \textbf{Serving:} Exposes all results through REST endpoints, persists to a serving database.
\end{enumerate}

\section{Container Architecture}
\label{sec:arch-containers}

The platform is decomposed into five containerised microservices, orchestrated by Docker Compose.

% ── Figure: Container Architecture ────────────────────────────────────────────
\begin{figure}[H]
\centering
\begin{tikzpicture}[node distance=1.5cm and 2cm]

  % External
  \node[external, minimum width=3cm] (user) {User / CVM\\REST Client};
  \node[external, minimum width=3cm, left=3cm of user] (ttdata) {TT Data\\(CSV / Excel)};
  
  % Services
  \node[service=ossblue, below=1.5cm of user, minimum width=3.5cm] (apigw) {
    \textbf{API Gateway}\\{\scriptsize FastAPI · :8000}
  };
  
  \node[service=aiviolet, below=of apigw, minimum width=3.5cm, xshift=-3.5cm] (ai) {
    \textbf{AI Service}\\{\scriptsize scikit-learn · :8001}
  };
  
  \node[service=cemteal, below=of apigw, minimum width=3.5cm, xshift=3.5cm] (worker) {
    \textbf{Pipeline Worker}\\{\scriptsize 5-phase orchestrator}
  };
  
  \node[datastore=datalake, below=1.5cm of ai, minimum width=3cm] (minio) {
    \textbf{MinIO}\\{\scriptsize Object Storage}\\{\scriptsize :9000}
  };
  
  \node[datastore=bssgreen, below=1.5cm of worker, minimum width=3cm] (pg) {
    \textbf{PostgreSQL 16}\\{\scriptsize Serving Store}\\{\scriptsize :5432}
  };
  
  % Docker boundary
  \begin{scope}[on background layer]
    \node[fit=(apigw)(ai)(worker)(minio)(pg),
          draw=darkbg!50, fill=darkbg!3, rounded corners=10pt,
          inner sep=18pt, line width=1pt, dashed] (docker) {};
  \end{scope}
  \node[below=2pt of docker.north east, anchor=north east,
        font=\scriptsize\bfseries, text=darkbg!60] {Docker Compose};
  
  % Arrows
  \draw[myarrow] (user) -- (apigw) node[arrlabel, midway, right] {REST};
  \draw[myarrow=cemteal!70] (ttdata) |- (worker) node[arrlabel, pos=0.75, above] {Ingest};
  
  \draw[myarrow] (apigw) -- (pg) node[arrlabel, midway, right, xshift=5pt] {Read};
  \draw[myarrow=cemteal!70] (worker) -- (ai) node[arrlabel, midway, above] {Inference};
  \draw[myarrow=cemteal!70] (worker) -- (minio) node[arrlabel, midway, left, xshift=-3pt] {Write};
  \draw[myarrow=cemteal!70] (worker) -- (pg) node[arrlabel, midway, right] {Persist};
  \draw[myarrow=aiviolet!70] (ai) -- (pg) node[arrlabel, midway, below, yshift=-2pt] {Results};
  \draw[myarrow=aiviolet!70] (ai) -- (minio) node[arrlabel, midway, below, yshift=-12pt] {Artifacts};
  
\end{tikzpicture}
\caption{Container-level architecture: five services orchestrated by Docker Compose.}
\label{fig:container-arch}
\end{figure}

\subsection{Service Descriptions}

\begin{table}[H]
\centering
\caption{Service inventory and responsibilities.}
\label{tab:services}
\begin{tabularx}{\textwidth}{@{}lllX@{}}
\toprule
\textbf{Service} & \textbf{Technology} & \textbf{Port} & \textbf{Responsibility} \\
\midrule
API Gateway  & FastAPI (Python)   & 8000 & REST endpoint serving: SLA risk, anomalies, correlations, revenue anomalies, pipeline runs \\
AI Service   & scikit-learn + FastAPI & 8001 & ML model hosting (GBR + IF $\times 2$), model training on startup, inference via POST \\
Pipeline Worker & Python (run-once) & --- & 5-phase orchestration: data ingestion, feature engineering, inference, decision \& action, curation \\
PostgreSQL 16 & Official image    & 5432 & Serving store with 8 tables: pipeline runs, datasets, models, anomalies, SLA risk, correlations, revenue anomalies, action recommendations \\
MinIO        & Official image     & 9000 & S3-compatible object storage: 3-layer data lake (raw/processed/curated) \\
\bottomrule
\end{tabularx}
\end{table}

\subsection{Inter-Service Communication}

All communication between services uses synchronous HTTP REST calls over the Docker internal network. The pipeline worker is the primary orchestrator --- it drives the entire data flow by calling the AI service for inference and writing results to both MinIO and PostgreSQL.

\section{Pipeline Architecture}
\label{sec:arch-pipeline}

The pipeline worker executes a 5-phase workflow that transforms raw data into actionable intelligence and autonomous recommendations.

% ── Figure: Pipeline Sequence ─────────────────────────────────────────────────
\begin{figure}[H]
\centering
\begin{tikzpicture}[
  step/.style={
    rectangle, rounded corners=2pt, minimum width=7cm, minimum height=0.6cm,
    draw=#1!60, fill=#1!8, font=\scriptsize, align=left, text=black,
  },
  phase/.style={
    font=\small\bfseries, text=#1!80!black, anchor=west,
  },
  node distance=0.15cm,
]

  % Phase 1: Ingestion
  \node[phase=cemteal] (p1) at (0,0) {Phase 1: Data Ingestion};
  \node[step=cemteal, below=0.3cm of p1, anchor=north west] (s1) at (0,-0.5) {
    1. Ensure MinIO buckets (raw / processed / curated)
  };
  \node[step=cemteal, below=of s1] (s2) {2. Ingest OSS data (real TT CSV/Excel or synthetic fallback)};
  \node[step=cemteal, below=of s2] (s3) {3. Ingest BSS data (real TT CSV/Excel or synthetic fallback)};
  \node[step=cemteal, below=of s3] (s4) {4--5. Upload raw data to MinIO raw layer};
  \node[step=cemteal, below=of s4] (s5) {6--7. Register pipeline run + raw datasets in PostgreSQL};
  
  % Phase 2: Processing
  \node[phase=ossblue, below=0.5cm of s5] (p2) {Phase 2: Feature Engineering};
  \node[step=ossblue, below=0.3cm of p2] (s6) {8--9. Process OSS $\to$ add severity, category, load factor, QoS score};
  \node[step=ossblue, below=of s6] (s7) {10. Process BSS $\to$ add ARPU category, data intensity, churn bucket};
  \node[step=ossblue, below=of s7] (s8) {11. Compute aggregate KPI features (9 OSS + 7 BSS metrics)};
  
  % Phase 3: AI Inference
  \node[phase=aiviolet, below=0.5cm of s8] (p3) {Phase 3: AI Inference};
  \node[step=aiviolet, below=0.3cm of p3] (s9) {12. POST /infer/sla-risk $\to$ GBR model $\to$ risk score [0, 1]};
  \node[step=aiviolet, below=of s9] (s10) {13. POST /infer/anomaly $\to$ IsolationForest (OSS)};
  \node[step=aiviolet, below=of s10] (s11) {14. POST /infer/revenue-anomaly $\to$ IsolationForest (BSS)};
  \node[step=aiviolet, below=of s11] (s12) {15. Compute OSS--BSS correlations (5 pairs $\times$ 2 methods)};
  
  % Phase 4: Decision & Action
  \node[phase=huaweired, below=0.5cm of s12] (p4) {Phase 4: Decision \& Action Engine};
  \node[step=huaweired, below=0.3cm of p4] (s13a) {Evaluate ML outputs against configurable business rules};
  \node[step=huaweired, below=of s13a] (s13b) {Generate action recommendations (alerts, churn prevention, upsell, escalation)};
  \node[step=huaweired, below=of s13b] (s13c) {Assign priority (critical/high/medium/low) and confidence score};
  
  % Phase 5: Persistence
  \node[phase=bssgreen, below=0.5cm of s13c] (p5) {Phase 5: Curation \& Persistence};
  \node[step=bssgreen, below=0.3cm of p5] (s14) {16--17. Build curated dataset (joined OSS+BSS+AI+Actions) $\to$ MinIO};
  \node[step=bssgreen, below=of s14] (s15) {Persist SLA risk, anomalies, revenue anomalies, models, correlations, action recommendations};
  
\end{tikzpicture}
\caption{The 5-phase pipeline execution flow.}
\label{fig:pipeline-flow}
\end{figure}

\section{Data Lake Architecture}
\label{sec:arch-datalake}

The platform implements a three-layer data lake using MinIO (locally) or OBS (on HCS):

% ── Figure: Data Lake Layers ──────────────────────────────────────────────────
\begin{figure}[H]
\centering
\begin{tikzpicture}[node distance=1.5cm]

  % Raw Layer
  \node[service=cemteal, minimum width=12cm, minimum height=1.8cm] (raw) {
    \textbf{Raw Layer}\\[2pt]
    {\scriptsize Real TT OSS data (cell KPIs) | Real TT BSS data (subscriber profiles)}\\
    {\scriptsize Anonymised at ingestion | Original format preserved | MinIO bucket: \texttt{raw}}
  };
  
  % Processed Layer
  \node[service=ossblue, minimum width=12cm, minimum height=1.8cm, below=of raw] (proc) {
    \textbf{Processed Layer}\\[2pt]
    {\scriptsize Feature-engineered OSS: severity, category, load factor, QoS score}\\
    {\scriptsize Feature-engineered BSS: ARPU category, data intensity, churn bucket, APPU, DOU}\\
    {\scriptsize MinIO bucket: \texttt{processed}}
  };
  
  % Curated Layer
  \node[service=aiviolet, minimum width=12cm, minimum height=1.8cm, below=of proc] (cur) {
    \textbf{Curated Layer}\\[2pt]
    {\scriptsize Joined OSS + BSS + AI outputs (anomaly flags, SLA risk, severity scores)}\\
    {\scriptsize Ready for analytics, dashboards, and CVM consumption | MinIO bucket: \texttt{curated}}
  };
  
  % AI Service on the side
  \node[service=aiviolet, right=1.5cm of proc, minimum width=2.5cm] (aiside) {
    \textbf{AI Service}\\{\scriptsize ML Inference}
  };
  
  % Arrows
  \draw[myarrow=cemteal!70, line width=1.2pt] (raw) -- (proc)
    node[arrlabel, midway, right, xshift=5pt] {Clean + Feature Engineer};
  \draw[myarrow=ossblue!70, line width=1.2pt] (proc) -- (cur)
    node[arrlabel, midway, right, xshift=5pt] {Join + Enrich with AI};
  \draw[myarrow=ossblue!70, line width=0.9pt] (proc) -- (aiside)
    node[arrlabel, midway, above] {Features};
  \draw[myarrow=aiviolet!70, line width=0.9pt] (aiside) |- (cur)
    node[arrlabel, pos=0.3, right] {Scores};
    
\end{tikzpicture}
\caption{Three-layer data lake architecture: Raw $\to$ Processed $\to$ Curated.}
\label{fig:datalake}
\end{figure}

\subsection{Layer Descriptions}

\begin{itemize}
  \item \textbf{Raw Layer:} Contains ingested data as received from Tunisie Telecom, after anonymisation (MSISDN hashing, address stripping). No transformations are applied. Data is stored as JSON/CSV objects in the \texttt{raw} MinIO bucket, preserving the original structure for auditability.
  
  \item \textbf{Processed Layer:} Contains feature-engineered data ready for ML inference. OSS records gain computed fields (latency severity, throughput category, load factor, QoS composite score). BSS records gain ARPU category, data intensity classes, churn risk buckets, and explicit APPU/DOU values. Stored in the \texttt{processed} bucket.
  
  \item \textbf{Curated Layer:} Contains the joined OSS+BSS dataset enriched with AI model outputs: anomaly flags, anomaly scores, SLA risk scores, and correlation insights. This is the ``gold'' layer, ready for dashboard consumption, CVM integration, or further analysis. Stored in the \texttt{curated} bucket.
\end{itemize}

\section{Database Architecture}

PostgreSQL 16 serves as the metadata and serving store. The schema comprises eight tables:

% ── Figure: ER Diagram (TikZ) ─────────────────────────────────────────────────
\begin{figure}[H]
\centering
\begin{tikzpicture}[
  entity/.style={
    rectangle, rounded corners=2pt, draw=#1!60, fill=#1!8,
    minimum width=3.8cm, font=\scriptsize, align=left,
    inner sep=5pt, line width=0.6pt,
  },
  entity/.default={darkbg},
  entitytitle/.style={font=\scriptsize\bfseries, text=#1!80!black},
  entitytitle/.default={darkbg},
  fk/.style={-{Stealth[length=4pt]}, draw=gray!60, line width=0.5pt},
  node distance=1.2cm and 2cm,
]

  % Pipeline Runs (center top)
  \node[entity=cemteal] (runs) {
    \textbf{pipeline\_runs}\\
    \rule{3.5cm}{0.3pt}\\
    id (PK, bigserial)\\
    run\_id (unique, text)\\
    status (text)\\
    started\_at (timestamptz)\\
    finished\_at (timestamptz)\\
    error\_message (text)
  };
  
  % Dataset Registry
  \node[entity=ossblue, below left=1.5cm and 0.5cm of runs] (datasets) {
    \textbf{dataset\_registry}\\
    \rule{3.5cm}{0.3pt}\\
    id (PK, bigserial)\\
    run\_id (FK)\\
    dataset\_type (text)\\
    layer (text)\\
    format (text)\\
    object\_key (text)\\
    schema\_version (text)\\
    row\_count (bigint)
  };
  
  % Model Registry
  \node[entity=aiviolet, below right=1.5cm and 0.5cm of runs] (models) {
    \textbf{model\_registry}\\
    \rule{3.5cm}{0.3pt}\\
    id (PK, bigserial)\\
    model\_name (text)\\
    version (text)\\
    artifact\_object\_key (text)\\
    created\_at (timestamptz)
  };
  
  % Anomalies
  \node[entity=huaweired, below=3cm of datasets, xshift=-1.5cm] (anomalies) {
    \textbf{anomalies}\\
    \rule{3.5cm}{0.3pt}\\
    id (PK, bigserial)\\
    run\_id (FK)\\
    cell\_id, region (text)\\
    kpi\_name (text)\\
    severity (float)\\
    value, baseline\_value (float)\\
    model\_name, model\_version
  };
  
  % SLA Risk Scores
  \node[entity=cvmorange, right=0.5cm of anomalies] (sla) {
    \textbf{sla\_risk\_scores}\\
    \rule{3.5cm}{0.3pt}\\
    id (PK, bigserial)\\
    run\_id (FK)\\
    region (text)\\
    window\_start, window\_end\\
    score (float, 0--1)\\
    explanation (JSONB)\\
    model\_name, model\_version
  };
  
  % Correlation Insights
  \node[entity=bssgreen, right=0.5cm of sla] (corr) {
    \textbf{correlation\_insights}\\
    \rule{3.5cm}{0.3pt}\\
    id (PK, bigserial)\\
    run\_id (FK)\\
    metric\_x, metric\_y (text)\\
    method (text)\\
    corr\_value (float)\\
    p\_value (float)
  };
  
  % Revenue Anomalies
  \node[entity=huaweired, right=0.5cm of corr] (revanom) {
    \textbf{revenue\_anomalies}\\
    \rule{3.5cm}{0.3pt}\\
    id (PK, bigserial)\\
    run\_id (FK)\\
    subscriber\_id (text)\\
    operator, line\_type (text)\\
    plan, metric\_name (text)\\
    severity (float)\\
    value, baseline\_value
  };
  
  % Foreign Key arrows
  \draw[fk] (datasets.north) -- (runs.south west);
  \draw[fk] (anomalies.north) -- (datasets.south);
  \draw[fk] (sla.north) -- (runs.south);
  \draw[fk] (corr.north) -- (runs.south east);
  \draw[fk] (revanom.north) -- (models.south east);
  
\end{tikzpicture}
\caption{Entity-Relationship diagram of the PostgreSQL serving store (8 tables).}
\label{fig:er-diagram}
\end{figure}

\section{API Architecture}

The API Gateway exposes seven RESTful endpoints:

\begin{table}[H]
\centering
\caption{REST API endpoints.}
\label{tab:api-endpoints}
\begin{tabularx}{\textwidth}{@{}llX@{}}
\toprule
\textbf{Method} & \textbf{Path} & \textbf{Description} \\
\midrule
GET & \texttt{/health}              & Service liveness check \\
GET & \texttt{/sla-risk}            & Latest SLA risk score from GBR model \\
GET & \texttt{/sla-risk/history}    & Last N SLA risk scores, newest first \\
GET & \texttt{/anomalies}           & Latest N OSS anomaly records with cell\_id, severity \\
GET & \texttt{/pipeline-runs}       & Last N pipeline execution records \\
GET & \texttt{/revenue-anomalies}   & Latest N BSS revenue anomalies \\
GET & \texttt{/correlation}         & Latest N OSS$\leftrightarrow$BSS correlations (Pearson/Spearman) \\
\bottomrule
\end{tabularx}
\end{table}

The AI Service exposes four internal endpoints (not externally accessible in production):

\begin{table}[H]
\centering
\caption{AI Service internal endpoints.}
\label{tab:ai-endpoints}
\begin{tabularx}{\textwidth}{@{}lllX@{}}
\toprule
\textbf{Method} & \textbf{Path} & \textbf{Model} & \textbf{Output} \\
\midrule
GET  & \texttt{/health}                & ---   & Model version, status \\
POST & \texttt{/infer/sla-risk}        & GBR & Risk score $\in [0, 1]$ + feature importances \\
POST & \texttt{/infer/anomaly}         & IF  & Per-record anomaly flag + normalised severity \\
POST & \texttt{/infer/revenue-anomaly} & IF  & Per-record revenue anomaly flag + severity \\
\bottomrule
\end{tabularx}
\end{table}

\section{Chapter Summary}

This chapter described the system architecture at four levels: CEM--CVM positioning, container decomposition, pipeline orchestration, and data lake structure. The platform comprises five Docker containers orchestrating a 5-phase pipeline that transforms raw telecom data into actionable intelligence and autonomous recommendations. The next chapter details the data strategy: how real data from Tunisie Telecom is ingested, anonymised, and processed.
